% Revised alternative to sections/2_related_v3.tex.
% Uses the existing aaai2026.bib and aaai_additions.bib files.
\section{Related Work}
\label{sec:related-work}

We draw on research on minority reception, minority expression and moral
reasoning, and interventions that shape exposure to different views.
These strands connect how dissent is received, how judgments and reasons
are expressed and discussed, and how readers respond when minority
threads are displayed.

\subsubsection*{Minority Reception}

Minority influence need not coincide with public acceptance
\cite{wood1994minority}. Dissenters can attract communication while also
facing rejection \cite{schachter1951deviation}. Attention to a minority
position therefore need not indicate approval of it.

Online feedback similarly separates evaluation from engagement. Prior
ratings can influence subsequent evaluations \cite{muchnik2013social},
and Reddit voting reflects local norms and practices
\cite{graham2021sociomateriality}. Negatively rated comments can
nevertheless attract more replies and larger reply trees
\cite{davis2021emotional}. In AITA, the relative valuation of agreement
and dissent also varies with the strength of the developing consensus
\cite{mellody2025consensus}. These studies motivate examining scores and
conversational engagement separately and considering the discussion's
level of agreement.
Reddit users who engage with opposing views receive fewer rewards in their political home communities \cite{Efstratiou2024}. We examine dissent within individual moral discussions, distinguishing approval from conversational engagement.

Reception may also differ with the moral direction of dissent. Public
blame can itself invite scrutiny and demands for justification
\cite{malle2022socially}. Blaming dissent assigns fault to a poster whom
the majority defends, whereas defending dissent challenges the
majority's blame. In \textbf{RQ1}, we examine how the relationship between
evaluation and conversational engagement varies with agreement strength
and the moral direction of dissent, using majority comments as a
reference.

\subsubsection*{Minority Discourse and Moral Reasoning}

AITA research maps everyday moral dilemmas and examines how disagreement
varies across situations and relationships
\cite{nguyen2022mapping,shi2025moral}. Within these discussions, judgments
express moral positions and contribute to exchanges of reasons.
Understanding minority discourse requires examining how firmly judgments
are expressed, how others respond, and what moral reasons emerge.

When people expect others to judge their views, they may try to please
their audience or defend what they already believe
\cite{tetlock1989strategies}. Research on the spiral of silence finds
that opinion climate is associated with expression among people with low
or moderate attitude certainty, but not those with high certainty
\cite{matthes2010spiral}. Anticipated personal attacks are also associated
with willingness to express minority views \cite{neubaum2018sanctions},
and expressions of doubt receive fewer recommendations in online news
discussions \cite{evans2023doubt}. These findings suggest that selective
participation may contribute to observed certainty patterns. Expressed
certainty must also be distinguished from private certainty.

Beyond how firmly judgments are expressed, replies can support or
challenge them. \citet{andreas2012annotating} annotate agreement and
disagreement in threaded discussions, interpreting responses in their
conversational context. We examine reply stance alongside moral
content, asking whether supportive and challenging replies introduce
reasons absent from their parent comments.

Computational research represents situation-specific moral norms
\cite{forbes2020social} and potentially conflicting values, rights, and
duties \cite{sorensen2024value}. \citet{xi2024morality} extract and
categorize reasons supporting moral judgments in AITA comments. These
approaches characterize moral content within contributions; examining
replies lets us identify how that content develops through interaction.
\textbf{RQ2} brings these dimensions together by examining expressed
certainty, reply stance, and moral reasons in comments and replies.

\subsubsection*{Information Exposure and Reader Responses}

Research on information exposure distinguishes available content,
algorithmic selection, and users' consumption choices
\cite{bakshy2015exposure}. Much intervention research concerns political
news and partisan attitudes, with varying effects of exposure to opposing
or like-minded sources
\cite{bail2018exposure,levy2021social,nyhan2023likeminded}.

Interventions differ in how they shape exposure. Social Mirror supports
reflection through network visualization and account recommendations
\cite{gillani2018mirror}, while Opinion Space and ConsiderIt introduce
dedicated interfaces for navigating perspectives or reasons
\cite{faridani2010opinion,kriplean2012considerit}. Other approaches change
which existing contributions appear prominently through comment
diversification \cite{northcutt2017diversification} or feed reranking
\cite{piccardi2025reranking}. Such selection changes what readers see
without adding a separate interface.

Closer experimental precedents compare agreeing and disagreeing comment
pairs of news \cite{dutceac2022cueing}.
Moral-judgment studies examine exposure to peer evaluations and arguments
\cite{kelly2017moral}, or aggregate controversy and group-certainty
signals derived from AITA discussions \cite{chen2026controversy}.
Our comparison combines comment selection with a moral-judgment task,
surfacing an existing minority thread about the same interpersonal
dilemma rather than an aggregate summary of community opinion.

Reader outcomes can extend beyond agreement. Dissent can encourage
divergent thinking \cite{nemeth1986differential}. Deliberative research
defines \emph{argument repertoire} as the reasons people can offer for
their own and opposing views \cite{cappella2002argument}, and associates
exposure to disagreement with a greater ability to offer opposing reasons
\cite{price2002disagreement}. These findings motivate examining readers'
consideration of reasons alongside their final judgments.

In \textbf{RQ3}, we test the replacement of a second majority thread with
the highest-scoring minority thread in a Reddit-style display, retaining
original wording, scores, and replies. Minority status is relative to the
discussion, rather than the reader's prior position. The experiment
compares these combined presentations of content and social cues. We
examine readers' judgment certainty and reported consideration of
alternative moral reasons, including among readers whose final judgment
agrees with the majority.
